\documentclass[10pt,a4paper]{amsart}
\usepackage[margin=19mm]{geometry}
\usepackage{amsmath,amssymb,mathtools}
\usepackage[scr=rsfs,cal=euler]{mathalfa}
\usepackage{xfrac}
\usepackage[unicode,hidelinks,bookmarks=false]{hyperref}
\newcommand{\ie}{i.e.\ }
\newcommand{\cf}{cf.\ }
\newcommand{\resp}{respectively\ }
\newcommand{\Subsection}{\subsection}
\providecommand{\nobreakdash}{\nobreak\hbox{-}}
\setlength{\emergencystretch}{2em}
\multlinegap0pt
\usepackage{xcolor}
\usepackage{amssymb}
\usepackage{enumerate}
\usepackage{mathtools}
\newtheorem{lemma}{Lemma}
\newcommand{\kint}{\vint}
\def\vint_#1{\mathchoice%
          {\mathop{\kern 0.2em\vrule width 0.6em height 0.69678ex depth -0.58065ex
                  \kern -0.8em \intop}\nolimits_{\kern -0.4em#1}}%
          {\mathop{\kern 0.1em\vrule width 0.5em height 0.69678ex depth -0.60387ex
                  \kern -0.6em \intop}\nolimits_{#1}}%
          {\mathop{\kern 0.1em\vrule width 0.5em height 0.69678ex
              depth -0.60387ex
                  \kern -0.6em \intop}\nolimits_{#1}}%
          {\mathop{\kern 0.1em\vrule width 0.5em height 0.69678ex depth -0.60387ex
                  \kern -0.6em \intop}\nolimits_{#1}}}
\title{Regularity for fully nonlinear elliptic equations in generalized Orlicz spaces}
\author{Sun-Sig Byun, Jeongmin Han, Mikyoung Lee}
\date{Source: arXiv:2512.16600v2 (CC BY 4.0)}
\begin{document}
\maketitle

\noindent\emph{Source and licence.} Regularity for fully nonlinear elliptic equations in generalized Orlicz spaces. Version arXiv:2512.16600v2 (CC BY 4.0), distributed by arXiv under the Creative Commons Attribution 4.0 International licence, http://creativecommons.org/licenses/by/4.0/, as recorded in the arXiv licence metadata for this exact version and shown on the abstract page https://arxiv.org/abs/2512.16600v2. Full licence notice: \texttt{licenses/arxiv-2512.16600-CC-BY-4.0.txt}.\par\smallskip

\noindent\textit{Editorial scope.} Counted unit: the article's lemma lemnon (source0.tex lines 909--929). The article's complete original proof of it is delivered with it (source0.tex lines 930--1002, one \texttt{proof} environment of 73 lines). No mathematics is written, repaired or re-derived: the statement and the proof are the article's own lines, in the article's own notation, and no retained line is substituted or re-flowed.  Retained uncounted as the source-local context the proof needs: lines 381--385; lines 615--648; lines 654--664; lines 830--839.  Nothing was omitted from any retained span, so the package carries no omission record.  No retained line contains a citation, so the wrapper carries no bibliography and no bibliography entry.  No retained line contains a figure environment, an image-inclusion command or drawing code, so no image is delivered and the package declares no asset dependency.  Editorial formatting only, in addition: an AMS \texttt{amsart} preamble that restates the article's own declarations and macros used by the retained lines, namely the environments \texttt{lemma} on separate counters, together with the standard packages those lines need.  The retained environments print in this document as Lemma 1 in order of appearance, whereas the article prints the same items with the numbering of its own full document.  The complete native source of the article as distributed in the arXiv e-print for this exact version is bundled unchanged as \texttt{sources/191258/source0.tex}.
\begin{lemma} \label{orljen}
Let $\varphi$ be a weak $\Phi$-function satisfying $\textrm{(aInc)}_{p}$ for some $p \ge 1$.
Then there exists $C=C(p,L)>0$ such that
$$ \varphi \bigg(C \bigg(\kint_{\Omega}|f|^{p} dx \bigg)^{\frac{1}{p}} \bigg)  \le \kint_{\Omega} \varphi(|f|)dx.$$
\end{lemma}

\begin{lemma}\label{arr_ac}
Assume that $F=F(X,x)$ is uniformly elliptic with constants $\lambda$ and $\Lambda$ with $F(0,\cdot)=0$,
and there exists $F^{\star}=F^{\star}(X,x)$ which is convex in $X$. Then the following hold:
\begin{itemize}
\item[(a)] Suppose that $f \in L^{\psi(\cdot)}(B_{6})$ and let $u$ be a viscosity solution to
\begin{align} \label{ob_muin2}
F_{\mu}(D^{2}u,x) = f \quad \textrm{in} \ B_{6}
\end{align}
with
$$ ||u||_{L^{\infty}(B_{6})} \le 1 \quad \textrm{and} \quad
\kint_{B_{6}}\psi(|f|,x) dx \le 1.$$
Then there exist small $\mu,\delta>0$ depending on $n,\lambda,\Lambda,p,q,L$ and $t_{0}$ such that if
$$ \bigg( \kint_{B_{r}(x_{0})\cap B_1} \sigma_{F^{\star}}(x_{0},x)^{n} dx  \bigg)^{\frac{1}{n}} \le \delta, $$
we have
$$ \kint_{B_{1}} \psi\big(\Theta (u,B_1),x\big) dx \le C$$
for some $C=C(n,\lambda,\Lambda,p,q,L,t_{0})>0$.
\item[(b)]
Suppose that $f \in L^{\psi(\cdot)}(B_{6}^{+})$ and let $u$ be a viscosity solution to
\begin{align} \label{ob_mubdry2}
\left\{ \begin{array}{ll}
F_{\mu}(D^{2}u,x) = f & \textrm{in $B_{6}^{+}, $}\\
u=0 & \textrm{on $T_{6}$}
\end{array} \right.
\end{align}
with
$$ ||u||_{L^{\infty}(B_{6}^{+})} \le 1 \quad \textrm{and} \quad
\kint_{B_{6}^{+}}\psi(|f|,x) dx \le 1.$$
Then there exist small $\mu,\delta>0$ depending on $n,\lambda,\Lambda,p,q,L$ and $t_0$ such that if  for any $r>0$ and $x_0\in \Omega$,
$$\bigg( \kint_{B_{r}(x_{0})\cap B_1^+} \sigma_{F^{\star}}(x_{0},x)^{n} dx  \bigg)^{\frac{1}{n}} \le \delta, $$
 we have
$$ \kint_{B_{1}^{+}} \psi\big(\Theta(u,B_{1}^{+}),x\big) dx \le C$$
for some $C=C(n,\lambda,\Lambda,p,q,L,t_{0})>0$.
\end{itemize}
\end{lemma}

\begin{lemma}\label{olbdest}
Let $\psi$ be a function depending only on $t$.
Then under the assumption of Lemma \ref{arr_ac},
there holds
%$$ \kint_{B_{2}} \psi(\Theta) dx \le C$$
%for some $C=C(n,\lambda,\Lambda,p,q,L,t_{0})>0$
%in the interior case, and
$$ \kint_{B_{2}^{+}} \psi\big(\Theta(u,B_{2}^{+})\big) dx \le C$$
for some constant $C=C(n,\lambda,\Lambda,p,q,L,t_{0})>0$.
%in the boundary case.
\end{lemma}

\begin{lemma} \label{stta}
Let $\lambda \ge \alpha \lambda_{0}$ with $\alpha = \big(\frac{120}{s_{2}-s_{1}} \big)^{n}$.
Then there exists a disjoint family $\{ B_{\tau_{k}}^{+}(y^{k}) \}_{k=1}^{\infty}$ with $y^{k} \in E(s_{1},\lambda)$ and $\tau_{k} \in (0, \frac{s_{2}-s_{1}}{60})$
such that
$$ E(s_{1},\lambda) \subset \bigcup_{k=1}^{\infty}B_{5\tau_{k}}^{+}(y^{k}),$$
and $\mathcal{H}_{y^{k}}(\tau_{k})=\lambda $ and $\mathcal{H}_{y^{k}}(\tau)<\lambda $ for any $\tau \in (\tau_{k},s_{2}-s_{1}]$,
where
$$\mathcal{H}_{y}(\tau) := \kint_{B_{\tau}^{+}(y)}\psi(\Theta, x)^{\frac{\nu}{nq}} dx+
\frac{1}{\delta} \bigg(\kint_{B_{\tau}^{+}(y)}\psi(|f|, x)^{\frac{1}{p}} dx\bigg)^{\frac{\nu p}{nq}} .$$
\end{lemma}

\begin{lemma} \label{lemnon}
Under assumptions in Lemma \ref{stta}, the following hold:
\begin{itemize}
\item[(a)] If $B_{30\tau_{k}}(y^{k}) \subset B_{s_{2}}^{+}$, we have
\begin{align*}
\bigg(\kint_{B_{30\tau_{k}}(y^{k})}\Theta^{\frac{\nu p}{q}} dx \bigg)^{^{\frac{q}{\nu p}}} \le C (\psi_{k}^{+})^{-1}\big(\lambda^{\frac{nq}{\nu}} \big)
\end{align*}and
\begin{align*}
\bigg(\kint_{B_{30\tau_{k}}(y^{k})}|f|^{n} dx \bigg)^{^{\frac{1}{n}}} \le C\delta^{\frac{q}{\nu p}} (\psi_{k}^{+})^{-1}\big(\lambda^{\frac{nq}{\nu}} \big)
\end{align*}
for some $C=C(n,\lambda,\Lambda,p,q,L,t_{0})>0$.
\item[(b)]  If $B_{30\tau_{k}}(y^{k}) \not\subset B_{s_{2}}^{+}$, we have
\begin{align*}
\bigg(\kint_{B_{60\tau_{k}}^{+}(\tilde{y}^{k}) }\Theta^{\frac{\nu p}{q}} dx \bigg)^{^{\frac{q}{\nu p}}} \le C (\psi_{k}^{+})^{-1}\big(\lambda^{\frac{nq}{\nu}} \big)
\end{align*}and
\begin{align*}
\bigg(\kint_{B_{60\tau_{k}}^{+}(\tilde{y}^{k}) }|f|^{n} dx \bigg)^{^{\frac{1}{n}}} \le C\delta^{\frac{q}{\nu p}} (\psi_{k}^{+})^{-1}\big(\lambda^{\frac{nq}{\nu}} \big)
\end{align*}
for some $C=C(n,\lambda,\Lambda,p,q ,L,t_{0})>0$, where $\tilde{y}^{k}=\big( (y^{k})',0 \big)$.
\end{itemize}
\end{lemma}

\begin{proof} 
We first consider the case $B_{30\tau_{k}} (y^{k})\subset B_{s_{2}}^{+}$.
Since $H_{y}^k(\tau) < \lambda$ for any $\tau \in [\tau_{k},  s_{2}-s_{1}]$ from Lemma \ref{stta},
we have
$$\kint_{B_{30\tau_{k}}(y^{k})}\psi(\Theta, x)^{\frac{\nu}{nq}} dx  \le \lambda
\quad \textrm{and} \quad \bigg(\kint_{B_{30\tau_{k}}(y^{k})}\psi(|f|, x)^{\frac{1}{p}} dx\bigg)^{\frac{\nu p}{nq}} \le \lambda\delta . $$

Setting $\theta (t)=\psi_{k}^{-}\big(t^{\frac{q}{\nu p}}\big)^{\frac{\nu}{nq}}$,
we see that $\theta$ satisfies $\textrm{(aInc)}_{1}$. Hence, by Lemma \ref{orljen},
we have
\begin{align*}
\psi_{k}^{-}\bigg( C \bigg( \kint_{B_{30\tau_{k}}(y^{k})} \Theta^{\frac{\nu p}{q}}dx
\bigg)^{\frac{q}{\nu p}} \bigg) &\le
\bigg( \kint_{B_{30\tau_{k}}(y^{k})} \psi_{k}^{-}(\Theta)^{\frac{\nu}{nq}} dx \bigg)^{\frac{nq}{\nu}}
\\ & \le \bigg( \kint_{B_{30\tau_{k}}(y^{k})} \psi(\Theta,x)^{\frac{\nu}{nq}} dx \bigg)^{\frac{nq}{\nu}} \\ & \le \lambda^{\frac{nq}{\nu}},
\end{align*}
where $C$ depends on $n$.

We also observe that
\begin{align*}
 \bigg( \kint_{B_{30\tau_{k}}(y^{k})}  \Theta^{\frac{\nu p}{q}}dx \bigg)^{\frac{nq}{\nu p}}
& \le \kint_{B_{30\tau_{k}}(y^{k})} \Theta^n dx
\\& \le C(\varphi^{-})^{-1}\bigg( \kint_{B_{30\tau_{k}}(y^{k})}
 \varphi^{-}( \Theta^n) dx \bigg) \\ & \le C_0(\varphi^{-})^{-1} \bigg(\frac{1}{|B_{5\tau_{k}}|} \bigg)
\end{align*} for some $C_0=C_0(n,\lambda,\Lambda,p,q,L,t_{0})>0$
from Lemma \ref{olbdest}.
Thus,
\begin{align*}
\varphi_{k}^{+} \bigg(\bigg( \kint_{B_{30\tau_{k}}(y^{k})} \Theta^{\frac{\nu p}{q}}dx \bigg)^{\frac{nq}{\nu p}} \bigg)&\le  C\bigg( \varphi_{k}^{+} \bigg( \frac{1}{C_0}\bigg( \kint_{B_{30\tau_{k}}(y^{k})} \Theta^{\frac{\nu p}{q}}dx \bigg)^{\frac{nq}{\nu p}}\bigg)  +1\bigg)
\\ & \le   C\bigg( \varphi_{k}^{-} \bigg(\bigg( \kint_{B_{30\tau_{k}}(y^{k})} \Theta^{\frac{\nu p}{q}}dx \bigg)^{\frac{nq}{\nu p}}\bigg)  +1\bigg)
\\ & = C\bigg( \psi_{k}^{-} \bigg(\bigg( \kint_{B_{30\tau_{k}}(y^{k})} \Theta^{\frac{\nu p}{q}}dx \bigg)^{\frac{q}{\nu p}}\bigg)  +1\bigg)
\\ & \le C\lambda^{\frac{nq}{\nu}},
\end{align*}
where $C=C(n,\lambda,\Lambda,p,q,L,t_{0})>0$, which implies
$$\bigg( \kint_{B_{30\tau_{k}}(y^{k})} \Theta^{\frac{\nu p}{q}}dx \bigg)^{\frac{q}{\nu p}}\le
C (\psi_k^{+})^{-1}(\lambda^{\frac{nq}{\nu}}).$$
Note that we have used the assumption (A1-$\varphi^{-}$) to derive the second inequality.


%Similarly as above, we can obtain the following estimates
%\begin{align*}
%C(n)\kint_{B_{5\tau_{k}}(y^{k})}|f|^{n} dx \le
%\big( (\psi^{-})^{-1}\big( \lambda\delta\big)^{\frac{nq}{\nu}}\big)^{n}
%\end{align*}
%and
%\begin{align*}
%\bigg(\kint_{B_{5\tau_{k}}(y^{k})}|f|^{n} dx \bigg)^{^{\frac{1}{n}}} \le
%C(\psi^{-})^{-1} \bigg(\frac{1}{|B_{5\tau_{k}}|} \bigg)
%\end{align*}
%for some $C=C(n,\lambda,\Lambda,p,q,L,t_{0})$.
Repeating a similar calculation, we finally obtain
\begin{align*}
\psi_{k}^{+}\bigg(\bigg(\kint_{B_{30\tau_{k}}(y^{k})}|f|^{n} dx \bigg)^{^{\frac{1}{n}}} \bigg) \le
C\big( \lambda\delta \big)^{\frac{nq}{\nu}}
\end{align*}
and this yields
$$\bigg(\kint_{B_{30\tau_{k}}(y^{k})}|f|^{n} dx \bigg)^{^{\frac{1}{n}}} \le C \delta^{\frac{q}{\nu p}}
(\psi_{k}^{+})^{-1}\big( \lambda^{\frac{nq}{\nu}} \big),$$
where $C=C(n,\lambda,\Lambda,p,q,L,t_{0})>0$. 

Next we consider the case $B_{30\tau_{k}}(y^{k}) \not\subset B_{s_{2}}^{+}$.
Observe that $|y^{k}-\tilde{y}^{k}|=|(y^{k})_{n}| < 30\tau_{k}$, $60\tau_{k} \le  s_{2}-s_{1} \le 1$ and
$y^{k} \in B_{s_{1}}^{+}$. Then we see that
$$B_{30\tau_{k}}^{+}(y^{k})\subset B_{60\tau_{k}}^{+}(\tilde{y}^{k}) \quad \textrm{and} \quad      B_{60\tau_{k}}^{+}(\tilde{y}^{k})\subset B_{s_1+60\tau_{k}}^{+} \subset  B_{s_{2}}^{+}.$$
Thus, by using a similar argument in the first case, we obtain
\begin{align*}
\bigg(\kint_{B_{60\tau_{k}}^{+}(\tilde{y}^{k}) }\Theta^{\frac{\nu p}{q}} dx \bigg)^{^{\frac{q}{\nu p}}} \le C (\psi_{k}^{+})^{-1}\big(\lambda^{\frac{nq}{\nu}} \big)
\end{align*}and
\begin{align*}
\bigg(\kint_{B_{60\tau_{k}}^{+}(\tilde{y}^{k}) }|f|^{n} dx \bigg)^{^{\frac{1}{n}}} \le C\delta^{\frac{q}{\nu p}} (\psi_{k}^{+})^{-1}\big(\lambda^{\frac{nq}{\nu}} \big)
\end{align*}
for some $C=C(n,\lambda,\Lambda,p,q,L,t_{0})>0$ by using the result of Lemma \ref{olbdest}.
\end{proof}

\end{document}
